\section{Conclusion}\label{sec:conclusion}

This paper studies startup subsidies in an environment in which entrepreneurial success often leads to acquisition rather than independent competition. Prospective entrepreneurs choose whether to enter, how much effort to exert, and whether to pursue an independence-oriented or a buyout-oriented project. After success, the entrant may be acquired by an incumbent, subject to merger approval, and a government chooses an entry subsidy before entry takes place.

The analysis yields three main results. First, a more permissive exit environment raises the private continuation value of entrepreneurship and therefore increases entry and effort. Second, it can redirect entrepreneurial project choice toward projects that are attractive primarily because they generate large acquisition premia rather than because they create strong stand-alone competitors. Third, because a uniform entry subsidy does not affect project choice, it cannot correct this composition distortion. Easier buyout exit need not strengthen the case for startup support and can instead reduce the optimal subsidy.

A central methodological implication is that welfare analysis must respect the same behavioral margins that govern private equilibrium. A merger is attempted only when expected private merger gains cover filing and implementation costs, so the social value of successful entry cannot be written as a simple convex combination of independent ownership and approved acquisition. In the paper's benchmark analysis, the resulting subsidy formulas and comparative statics are derived first from product-market primitives and the endogenous merger-attempt condition alone. Once this distinction is imposed, merger policy affects subsidy design only when it changes actual behavior.

The paper also characterizes a local--national incidence wedge. Local and national authorities may value successful entrepreneurial projects differently because the incidence of gains and losses from acquisition differs across levels of government. In the revised presentation, those local and national results are derived as planner-specific incidence adjustments around the primitive benchmark. For a fixed project, the local--national difference is therefore an incidence comparison. What buyout exit adds is that the privately selected project can itself change with merger permissiveness, so the economically relevant local--national comparison can change as well. Local governments may therefore subsidize too much or too little relative to the national benchmark even when both policymakers face the same private equilibrium.

The broader implication is that startup subsidies and merger policy should not be studied in isolation. Entry subsidies govern the scale of entrepreneurship, but the exit environment shapes private project direction. When acquisition is an important exit route, policy design requires instruments that address both margins. A final, explicitly exploratory implication concerns policy design beyond the baseline model: buyout-contingent clauses, such as clawbacks tied to acquisition, can complement uniform subsidies by affecting not only how much entrepreneurship occurs but also what type of entrepreneurship is pursued. This extension is intended to illustrate instrument separation rather than to displace the paper's main analytical contribution, which lies in the behaviorally consistent welfare analysis of Sections \ref{sec:subsidy} and \ref{sec:decentralized}.

The model is intentionally stylized. Natural extensions include richer merger review, multiple entrants, dynamic innovation races, and explicit multi-jurisdiction competition. The present paper, however, studies incidence-based divergence between local and national policy rather than strategic interaction across jurisdictions. Even so, the main message is simple: in economies where acquisition is a central entrepreneurial exit route, startup policy is not only about creating more firms. It is also about shaping what kinds of firms are created and how the gains and losses from their exit are distributed across policymakers.
